\section{A regular point in the trace family}\label{sec:regular-point}

The admissible module of Proposition~\ref{prop:diagonal} still need not
have an algebraic vector.  We first show that the size of its commuting
trace family forces positive Iwasawa rank.  That rank supplies algebraic
vectors with a prescribed smooth type.  A coefficient eigensystem among
these vectors will give the regular point needed in the final tangent
argument.

Keep the domain $D$, local block ring $R$, and nonzero finite
$D$-torsion-free module $L$ of Proposition~\ref{prop:family}.  Thus
$\dim D\geq3$, $D$ is finite over $R$, and
\[
 V=L\widehat\otimes_{\mathcal E}P
\]
is $\OO$-torsion-free and finite over $\mathscr A=\OO[[K]]$.
Here $K$ is a sufficiently deep uniform subgroup of $\SL_2(\Z_2)$.
All coefficient fields below are finite extensions of $\Q_2$.

\subsection{Dimension forces positive rank}

Write $\mathfrak a=(\varpi,\,k-1:k\in K)$ for the maximal
augmentation ideal of $\mathscr A$.  We use its filtration throughout.
After shrinking $K$, the ordered-monomial description of a uniform
Iwasawa algebra gives
\begin{equation}\label{eq:iwasawa-graded}
 \operatorname{gr}_{\mathfrak a}\mathscr A
       \simeq\kk[X_0,X_1,X_2,X_3].
\end{equation}
For $\Z_2$ coefficients this is
\cite[Theorem~3.22 and Lemmas~3.24--3.25]{Venjakob2002}.
The same statement for $\OO$ follows from the ordered expansion in
$\varpi$ and $b_i=k_i-1$, for three ordered uniform generators $k_i$:
take $K$ deep enough that its commutator relations have
$\mathfrak a$-degree at least three.  The symbols of
$\varpi^{a_0}b_1^{a_1}b_2^{a_2}b_3^{a_3}$ then give the polynomial
basis in~\eqref{eq:iwasawa-graded}.

For a nonzero finite $\mathscr A$-module $M$, let $g(M)$ be the degree
of the eventual Hilbert polynomial
\[
 H_M(n)=\operatorname{length}_{\OO}(M/\mathfrak a^nM).
\]
Good filtrations, the Artin--Rees property, and
\eqref{eq:iwasawa-graded} imply $0\leq g(M)\leq4$.
These are the filtered-algebra facts used also in the proof of
\cite[Proposition~4.4]{FM2}.

The decisive comparison is between the dimension of $D$ and this
growth degree. We will make $\dim D-1$ successive prime specializations
to a coefficient fiber chosen to have growth degree at least two:
its reduction modulo $\varpi$ is infinite-dimensional, and $\varpi$
acts injectively. The following lemma shows that each specialization
lowers growth by at least one. Thus $\dim D\geq3$ forces the maximal
growth degree four.

\begin{lemma}\label{lem:growth-drop}
Let $M$ be a finite $\mathscr A$-module and let $f:M\to M$ be
injective and $\mathscr A$-linear.  If $M/fM\ne0$, then
\[
                 g(M/fM)\leq g(M)-1.
\]
Moreover, a finite $\mathscr A$-module of growth degree four has
positive rank over $\mathscr A$.
\end{lemma}
\begin{proof}
The induced filtration on $fM$ is good.  Artin--Rees, followed by the
isomorphism $M\simeq fM$, therefore gives an integer $c\geq1$ such that
\[
 fM\cap\mathfrak a^nM\subseteq f(\mathfrak a^{n-c}M)
 \qquad(n\geq c).
\]
Taking lengths in the induced exact sequence yields
\[
 H_{M/fM}(n)
 \leq H_M(n)-H_M(n-c)=O(n^{g(M)-1}).
\]
This proves the first assertion, including the impossibility of a
nonzero cokernel when $g(M)=0$.

The algebra $\mathscr A$ is a noetherian domain and has an Ore
division ring of fractions.  If a finite module has rank zero,
each element is annihilated by a nonzero element of $\mathscr A$.
A filtration by its finitely many generators has cyclic torsion
quotients.  Such a quotient is of the form $\mathscr A/I$, where
$I$ contains a nonzero element $a$.  Its associated graded module
is a quotient of the polynomial ring in~\eqref{eq:iwasawa-graded}
by the nonzero initial form of $a$, and hence has growth degree at
most three.  The same bound holds for a finite extension of these
quotients.  Thus growth degree four implies positive rank.
\end{proof}

\begin{proposition}\label{prop:positive-rank}
The module $V$ satisfies
\[
                       \rank_{\mathscr A}V>0.
\]
\end{proposition}
\begin{proof}
Choose an auxiliary coefficient point $y$ of $D$ whose local
semisimple parameter, if reducible, has character ratio outside
$\{1,\eps,\eps^{-1}\}$.  Such points exist.  Indeed, at a reducible
exceptional parameter the fixed determinant and the specified ratio
determine each character up to a quadratic character.  Local class
field theory gives only finitely many quadratic characters of
$\Q_2^\times$.  The exceptional parameters therefore lie in a finite
union of coefficient-point closures in $\Spec R$, of integral
dimension at most one.  Finiteness of $D$ over $R$ gives the same
bound for their inverse image.  Since $D$ is horizontal of dimension
at least three, coefficient-point density
\cite[Lemma~3.3]{FM2} supplies $y$ outside that inverse image.

Put $s=\dim D$.  If $\mathfrak p_y$ is the kernel of the coefficient
point, then $D/\mathfrak p_y$ is finite over $\OO$ and has dimension
one.  The complete local domain $D$ is catenary, so there is a
saturated chain
\[
 (0)=\mathfrak p_0\subsetneq\mathfrak p_1
       \subsetneq\cdots\subsetneq\mathfrak p_{s-1}=\mathfrak p_y.
\]
Starting with $L^{(0)}=L$, define $L^{(i+1)}$ by tensoring $L^{(i)}$ with
$D/\mathfrak p_{i+1}$ and removing its torsion over that domain.
Each $L^{(i)}$ is nonzero: the preceding torsion-free finite module has
full support, and Nakayama's lemma at the next prime shows that its
fiber there is nonzero.  Thus $L^{(i+1)}$ has nonzero generic fiber.
The commuting $\mathcal E$-action survives all these operations.

Choose $a_i\in\mathfrak p_{i+1}\setminus\mathfrak p_i$.
Multiplication by $a_i$ is injective on $L^{(i)}$.  The exact compact
block equivalence of Theorem~\ref{input:local-blocks} makes it injective
on
\[
                  V_i=L^{(i)}\widehat\otimes_{\mathcal E}P,
\]
and $V_{i+1}$ is a quotient of its cokernel.  All these objects are
finite over $\mathscr A$, being successive quotients of $V$.
Lemma~\ref{lem:growth-drop} gives
\begin{equation}\label{eq:chain-growth}
             g(V)\geq s-1+g(V_{s-1}).
\end{equation}

Finally, $L^{(s-1)}$ is nonzero and finite torsion-free over $\OO$.
Its associated Banach representation has finite length and, after
a finite coefficient extension, every simple factor has the
central parameter of a coefficient conjugate of $y$; see
Theorem~\ref{input:local-blocks}.
They are infinite-dimensional: the generic principal-series
possibilities are infinite-dimensional, as is the representation
attached to an absolutely irreducible two-dimensional Galois
parameter.  Our choice of $y$ excludes the exceptional character
possibilities.  Hence $V_{s-1}/\varpi V_{s-1}$ is infinite-dimensional
over $\kk$.  Otherwise compact Nakayama would make $V_{s-1}$ finite
over $\OO$, contradicting this Banach description.

Exactness also makes multiplication by $\varpi$ injective on
$V_{s-1}$.  The infinite-dimensional reduction has growth degree at
least one, so Lemma~\ref{lem:growth-drop} gives $g(V_{s-1})\geq2$.
Together with~\eqref{eq:chain-growth}, this proves
\[
                    4\geq g(V)\geq\dim D+1\geq4.
\]
The final assertion of Lemma~\ref{lem:growth-drop} finishes the proof.
\end{proof}

\subsection{Capturing a fixed smooth type}

Let $K_0=\GL_2(\Z_2)$ and let
\[
 \sigma:K_0\longrightarrow\GL_2(\F_2)\simeq S_3
              \xrightarrow{\operatorname{sign}}\{1,-1\}
\]
be the inflated sign character over $E$. Its restriction to the upper
unipotent subgroup of $\GL_2(\F_2)$ is nontrivial, so it is a cuspidal
representation of that finite group. We use this type for two reasons.
Tensoring it with an algebraic representation excludes every Banach
simple with reducible Galois parameter. In a classical smooth representation it
forces depth-zero supercuspidality, and hence a Weil--Deligne parameter
trivial on wild inertia. The second property will exclude the CM fields
inside $\Q(\zeta_{2^\infty})$ and thereby allow the adjoint Selmer
vanishing used in Proposition~\ref{prop:tangent-bound}.

Write
\[
          \Pi_V=\Hom_{\OO}^{\cts}(V,\OO)[1/2]
\]
for the unitary admissible Banach representation associated with $V$.
Its central character on scalar units is $z\mapsto z^{d}$, where
$d=w+1$ by the determinant normalization.  Negative determinant
powers are permitted in the algebraic representations below.

\begin{lemma}\label{lem:type-capture}
There is an irreducible algebraic representation
$U=\Sym^n(E^2)\otimes\det^a$, with $n\geq0$, $a\in\Z$, and
$n+2a=d$, such that
\[
                  \Hom_{K_0}(\sigma\otimes U,\Pi_V)\ne0.
\]
This Hom space is finite-dimensional and carries a continuous
commuting action of $D$.
\end{lemma}
\begin{proof}
Positive rank supplies a nonzero $\mathscr A$-linear map
$V\to\mathscr A$: take a nonzero linear functional after passing
to the Ore division ring and clear right denominators on finitely
many generators.  Maps between finite $\mathscr A$-modules are
continuous.  Choose $K$ so that its intersection with the scalar
unit subgroup $Z_0$ is trivial.  Extend the action on the target
$\mathscr A$ to $Z_0K$ by the dual central character $z^{-d}$.
Since $Z_0K$ has finite index in $K_0$, Frobenius reciprocity gives
a nonzero map from $V$ to its coinduced compact module.  Dualizing
gives a nonzero continuous $K_0$-map
\begin{equation}\label{eq:functions-to-banach}
 \mathcal C_d\longrightarrow\Pi_V,
 \qquad
 \mathcal C_d=\{f\in C(K_0,E):f(kz)=z^{-d}f(k)\}.
\end{equation}
The action on $\mathcal C_d$ is left translation.

We verify the density needed to use this map.  Matrix coefficients
of algebraic representations of central degree zero form an
algebra separating the points of $K_0/Z_0$: coefficients of the
adjoint representation already separate them.  Polynomial
approximation on compact subsets of a finite-dimensional
$2$-adic affine space therefore makes this algebra dense in
$C(K_0/Z_0,E)$.  Concretely, embed the compact quotient by the
adjoint matrix entries and apply multivariable polynomial
approximation on a containing compact box, after extending
continuous functions by clopen partitions.

Fix one algebraic representation of central degree $d$, for example
$\det^{d/2}$ if $d$ is even and $E^2\otimes\det^{(d-1)/2}$ if $d$
is odd.  Its matrix coefficients have no common zero on $K_0$.
On a finite clopen cover of $K_0/Z_0$, divide a function in
$\mathcal C_d$ by a nonvanishing coefficient of this representation
evaluated at $k^{-1}$.  Approximate the resulting functions of central
degree zero as above, and multiply back.  Complete reducibility of
algebraic representations shows that the resulting dense span consists
of matrix coefficients of irreducible algebraic representations with
central degree $d$.  Multiplication by the nowhere-zero function
$\sigma(k^{-1})$ preserves $\mathcal C_d$ and its topology.  Thus
the matrix coefficients of the types $\sigma\otimes U$ in the
statement also have dense span.  A nonzero continuous map in
\eqref{eq:functions-to-banach} cannot kill them all.

For finiteness, transpose a type map to a continuous $K$-equivariant
map from $V$ to the dual of that finite-dimensional type.  A finite
set of $\mathscr A$-generators of $V$ determines such a map, so the
space of maps is finite-dimensional.  Evaluation on those generators
also proves continuity of its commuting $D$-action.
\end{proof}

\begin{lemma}\label{lem:type-irreducible}
Let $E'/E$ be a finite coefficient extension, let $\Pi$ be one of the
$E'$-Banach simples described in Theorem~\ref{input:local-blocks}, and
let $U=\Sym^n(E'^2)\otimes\det^a$, with $n\geq0$ and $a\in\Z$.
If
\[
                  \Hom_{K_0}(\sigma\otimes U,\Pi)\ne0,
\]
then the Galois parameter of $\Pi$ is absolutely irreducible.
\end{lemma}
\begin{proof}
Suppose the parameter is reducible. Theorem~\ref{input:local-blocks}
lists its possible Banach simples. Put
$S_0=\SL_2(\Z_2)$ and
$u=\left(\begin{smallmatrix}1&1\\0&1\end{smallmatrix}\right)$.
On $\sigma\otimes U$, the operator $u$ is minus a unipotent
operator.  In particular $u-1$ is invertible over $E'$.
The group $S_0$ acts transitively on $\mathbb P^1(\Q_2)$.
An $S_0$-map from the type $\sigma\otimes U$ to a continuous principal series, followed
by evaluation at the identity coset, would give a functional invariant
under the upper unipotent subgroup, hence under $u$.  This functional
is zero, and transitivity makes the whole map zero.  The same argument
excludes a character, whose restriction to $S_0$ is trivial.

For a twist of the continuous Steinberg representation, pull back
the defining quotient of a continuous induction by a character along
a proposed type map.  This gives an extension of our finite-dimensional
type by a character.  Every finite-dimensional continuous
representation of $S_0$ over $E'$ is semisimple.  To see this,
continuity and logarithm charts make it analytic on a sufficiently
small open subgroup.  Complete reducibility for
$\mathfrak{sl}_2(E')$ supplies an equivariant projection onto any
invariant subspace.  That projection commutes with a sufficiently
small open normal subgroup of $S_0$; averaging its conjugates over
the finite quotient gives an $S_0$-equivariant projection.
The pulled-back extension therefore splits over $S_0$, producing
the already excluded map into continuous induction.

This excludes every reducible parameter in the classification.
\end{proof}

\subsection{From a type eigenvector to a classical point}

\begin{proposition}\label{prop:regular-point}
There is a coefficient point $x:D\to\OO'$ at which the global
pseudorepresentation is that of a continuous absolutely irreducible
representation
\[
                    r_x:G_{\Q}\longrightarrow\GL_2(E'),
                    \qquad E'=\Frac(\OO'),
\]
after a finite coefficient extension if necessary.  This representation
is odd and finitely ramified.  Its restriction to $G_{\Q_2}$ is
absolutely irreducible and regular de Rham, and
$\WD(r_x|_{G_{\Q_2}})$ is trivial on wild inertia.  Consequently
$r_x$ is attached to a classical cuspidal eigenform up to Tate twist.
\end{proposition}
\begin{proof}
Take a type supplied by Lemma~\ref{lem:type-capture}.  The commuting
$D$-action on its nonzero finite-dimensional Hom space has a common
eigenvector after finite coefficient extension.  It gives a continuous
character $x:D\to E'$.  Since $D$ is compact, its image in the
endomorphism algebra preserves a lattice; its eigenvalues are integral.
Thus $x$ takes values in $\OO'$, and the corresponding map
\begin{equation}\label{eq:type-eigenspace}
                   \sigma\otimes U\longrightarrow\Pi_V[x]
\end{equation}
is nonzero.

This eigenspace has finite length as a Banach representation.
Indeed, extend coefficients to $\OO'$, impose the scalar relations
$a=x(a)$ for $a\in D$ on $L$, and remove $\OO'$-torsion.  The
result is finite over $\OO'$.  The exact compact equivalence identifies
its inverse with the corresponding torsion-free scalar quotient of
$V$; continuous maps to $\OO'$ kill the discarded torsion.
Dualizing and inverting $2$ therefore identifies its Banach
representation with $\Pi_V[x]$.  The finite-length assertion is the
coefficient-fiber assertion of Theorem~\ref{input:local-blocks}.
A composition series now shows that some Banach simple contains
the type in~\eqref{eq:type-eigenspace}: follow a nonzero map through
the series, passing to a quotient when its composite is nonzero and
otherwise factoring it through the closed subrepresentation.
No exactness assertion for the functor of locally algebraic vectors
is involved.

Lemma~\ref{lem:type-irreducible} now makes the local block parameter
$q_x$ absolutely irreducible.  The global determinant at $x$ is then absolutely
irreducible as well and is represented, after finite coefficient
extension, by a continuous $r_x$; this is the representation theorem
for characteristic-zero determinants
\cite[Theorem~2.12]{Chenevier2014}.
Continuity can also be checked by recovering matrix coordinates
from traces against a matrix-algebra basis.
The normalization is
\begin{equation}\label{eq:regular-local-normalization}
                   q_x=r_x|_{G_{\Q_2}}\otimes\eps.
\end{equation}
Since $2$ splits completely in $F$, restriction to $G_F$ is still
absolutely irreducible.

We can now transport the type to the classical forms upstairs.
The point $x$ induces a Hecke coefficient point through
$\mathbb T\to D$.  By \cite[Proposition~4.2]{FM2}, its completed
definite-forms eigenspace contains a Banach tensor product
\[
                         \widehat\bigotimes_{v\mid2}\Pi_v,
\]
where $\Pi_v$ is a simple with parameter $q_x$.  All these parameters
are identical, and an absolutely irreducible parameter has a unique
Banach simple.  Thus each $\Pi_v$ contains the type already found,
and their tensor contains
\[
                         \bigotimes_{v\mid2}(\sigma\otimes U).
\]
These are locally algebraic vectors in the eigenspace of completed
forms.  The classical comparison
\cite[Proposition~4.1]{FM2} identifies them with definite automorphic
forms of algebraic factor $\bigotimes U$.  The weights are regular:
the factor $\Sym^n$ gives classical weight $n+2\geq2$, while the
determinant power only changes the common twist.  In the classical
decomposition the smooth factor contains $\sigma$ at every dyadic
place.  Indeed the Lie algebra identifies the algebraic factor,
and the remaining $K_0$-action is its smooth multiplicity space.
Reduced-norm characters cannot contain this type.  Jacquet--Langlands
therefore gives a regular cuspidal Hilbert eigensystem over $F$.
Its Galois representation agrees with $r_x|_{G_F}$ by the Hecke
identities and Chebotarev.

The finite-group cuspidality of $\sigma$, established above, now
identifies the local smooth factors. The depth-zero type theorem
identifies its characteristic-zero smooth occurrences with depth-zero
supercuspidal representations
\cite[Appendix~A, Sections~A.3.1--A.3.2]{Henniart2002}.
The local Langlands correspondence preserves depth
\cite[Theorem~2.9]{ABPS2016}; its corresponding Weil parameter
therefore has depth zero, which means that it is trivial on wild
inertia.
The dyadic local--global comparison in
\cite[Proposition~4.1]{FM2} now gives regular de Rhamness and this
wild-inertia assertion for $r_x|_{G_F}$ at every dyadic place.
The splitting of $2$ in $F$ gives exactly those properties for
$r_x|_{G_{\Q_2}}$.  The cyclotomic twist in
\eqref{eq:regular-local-normalization} has unramified Weil--Deligne
character and does not change the wild-inertia assertion.

Finally, finite ramification is built into $D$, and its fixed
determinant $\chi$ is odd.  All the hypotheses of
\cite[Theorem~1.1]{FM2} have now been verified for $r_x$: continuity,
absolute irreducibility, oddness, finite ramification, and two
distinct de Rham weights at $2$.  That theorem gives the asserted
classical modularity up to Tate twist, with no restriction on the
residual representation.
\end{proof}
